\section{Institutional Setting and Research Question}
\label{sec:setting}

In February 2024, the Securities and Exchange Commission shortened the initial
Schedule~13D filing deadline from ten calendar days to five business days, the
first change to the accumulation window since 1968
\parencite{sec2023modernization}. Once an investor crosses five per cent
beneficial ownership, the amendment reduces the period in which further shares
may be acquired before the position is disclosed. The final rule also declined
to adopt proposed Rule~13d-3(e), which would have treated certain cash-settled
derivatives as beneficial ownership. The reform therefore changed both the
period for undisclosed equity accumulation and the relative attractiveness of
derivative exposure.

The paper asks whether this change altered activist campaign formation, stake
formation, and the distribution of takeover gains. It addresses three
questions. First, did campaign entry or target composition change more for
targets on which the deadline was more restrictive? Second, conditional on a
campaign, did reported stakes and the use of cash-settled derivatives change
with that exposure? Third, did activist-related takeover gains appear at filing,
during the pre-bid run-up, or in the announced bid markup?

Exposure is measured by the required trading days,
\begin{equation}
\mathrm{RTD}_i
=\frac{0.05\times S_{\mathrm{out},i}}{\phi\times\mathrm{ADV}_i},
\label{eq:rtd-intro}
\end{equation}
the number of trading days required to acquire an additional five percentage
points of ownership in target $i$ at participation cap $\phi$. Here
$S_{\mathrm{out},i}$ denotes shares
outstanding and $\mathrm{ADV}_i$ is average daily volume, measured in a lagged
pre-trigger window. Since the statutory clock begins only at the five-per-cent
crossing, the deadline limits the undisclosed increment beyond that threshold
to $0.05\times d/\mathrm{RTD}_i$ percentage points. Higher
$\mathrm{RTD}_i$ therefore indicates a tighter constraint for a given legal
window $d$, without using the campaign's realised accumulation or outcome.

\section{Economic Mechanism and Predictions}
\label{sec:predictions}

A shorter filing window affects activist activity through two margins. It
raises the price-impact cost of acquiring a given undisclosed increment, and it
caps the increment that can be acquired before disclosure. These effects may
discourage campaigns or shift them towards more liquid targets, larger
activists, or different campaign objectives. They may also encourage
substitution from common-stock accumulation to cash-settled derivatives reported
in Item~6, an instrument that proposed Rule~13d-3(e) would have covered but the
final rule did not.

The economic magnitude is uncertain. The SEC projected approximately
\$810 million a year in foregone shareholder value, reasoning that less time for
undisclosed accumulation would deter value-increasing campaigns
\parencite{sec2023modernization}; \textcite{polk2024} project lower activist
returns using historical data. At the same time, the SEC's tables indicate that
roughly 80 per cent of filers had already completed accumulation by the
five-business-day deadline and that the rule binds materially for only a
minority of campaigns. The central empirical question is consequently not
whether the reform must chill activism, but whether any response is concentrated
among targets for which the deadline plausibly binds.

The interpretation of takeover premia is likewise not mechanical.
\textcite{celentanolevine2025} find lower conditional bid premia for activist
targets and relate this result to bargaining. Yet the pre-bid price may already
reflect information revealed at filing and during subsequent trading.
\textcite{schwert1996} finds little substitution between run-ups and markups,
whereas \textcite{bett2014} finds stronger substitution after adjusting the
run-up for changes in stand-alone value. I therefore decompose total target gain
into the filing return $G_{\mathrm{fil}}=\log(P^F/P^0)$, the post-filing run-up
$G_{\mathrm{run}}=\log(P^B/P^F)$, and the bid markup
$G_{\mathrm{mkp}}=\log(B/P^B)$. The decomposition is evaluated over a fixed
trigger-to-bid horizon so that variation in event timing cannot mechanically
generate the predicted pattern.

\section{Contribution}
\label{sec:contribution}

The contribution is a realised evaluation of the February 2024 Schedule~13D
acceleration using predetermined variation in required trading days. Existing
disclosure-threshold models study stake-size triggers; models of activist
intermediation and bidder solicitation take the toehold as given or make
disclosure instantaneous; and structural work does not vary the filing window.
The model here makes the legal deadline a time constraint on post-threshold
accumulation and yields the exposure measure used in the empirical design. The
paper also tests whether the reform affected derivative use and records the
timing of activist-related takeover gains. It does not claim to re-establish the
activist-as-intermediary mechanism or the capitalisation interpretation of
conditional premia. Those results instead motivate the empirical tests.
